%% ma: L11-B17-0021
%% lop: 11
%% chuong: 5
%% bai: 17
%% dang: 11.17.4
%% loai: TN4
%% muc: NB
%% dap_an: "C"
%% nguon: "tai-lieu-hoc-tap-mon-toan-lop-11-theo-chuong-trinh-giao-duc-pho-thong-2018.pdf (Lớp Toán Thầy Duy, Phạm Lê Duy, THPT Chu Văn An), Chương 3 Bài 3, Phần 1, Câu 1 (tr.398)"
%% kien_thuc: "Lớp 11 Bài 17 (hàm số liên tục)"
%% trang_thai: cho-duyet
%% ly_do: "Dạng mới đề xuất (11.17.x), chờ giáo viên duyệt dạng; tài liệu gốc không có lời giải mẫu, đáp án do kho tự giải và kiểm bằng sympy."
%% ghi_chu: "Đề gốc in 'liên tục trên cm' (lỗi chữ, mất ký hiệu); kho khôi phục thành 'liên tục trên khoảng (a;b)' theo định nghĩa 2 trong lý thuyết, đáp án gốc C khớp."
\begin{ex}
Cho hàm số $y=f(x)$ liên tục trên khoảng $(a;b)$. Điều kiện cần và đủ để hàm số liên tục trên $[a;b]$ là
\choice{$\lim\limits_{x\to a^+}f(x)=f(a)$ và $\lim\limits_{x\to b^+}f(x)=f(b)$.}{$\lim\limits_{x\to a^-}f(x)=f(a)$ và $\lim\limits_{x\to b^-}f(x)=f(b)$.}{\True $\lim\limits_{x\to a^+}f(x)=f(a)$ và $\lim\limits_{x\to b^-}f(x)=f(b)$.}{$\lim\limits_{x\to a^-}f(x)=f(a)$ và $\lim\limits_{x\to b^+}f(x)=f(b)$.}
\loigiai{\sloppy Theo định nghĩa, hàm số liên tục trên $(a;b)$ liên tục trên đoạn $[a;b]$ khi $\lim\limits_{x\to a^+}f(x)=f(a)$ và $\lim\limits_{x\to b^-}f(x)=f(b)$ (chỉ xét giới hạn phía nằm trong $[a;b]$).}
\end{ex}
